\section*{الفصل \ref{ch:b2:isotope-geochemistry} --- الجيوكيمياء النظائرية}

\begin{solution}{exo:b2:isotope-geochemistry:1}
تريستان: $(0.51249/0.512638 - 1)\times10^4 = -2.9$. مانغايا: $(0.51285/0.512638 -
1)\times10^4 = +4.1$. بازلت MORB: $+8.5$. فمصدرا مانغايا وMORB، وقيمة
$\varepsilon_{\mathrm{Nd}}$ فيهما موجبة، كانت لهما في المتوسط نسبة Sm/Nd أعلى من الكوندريتات؛
وكانت لمصدر تريستان نسبة أدنى.
\end{solution}

\begin{solution}{exo:b2:isotope-geochemistry:2}
$\delta^{18}\mathrm{O}_{\mathrm{VPDB}} = (\num{2.0672e-3}/\num{2.0052e-3} - 1)\times1000 =
+30.9$ بالجزء في الألف على مقياس \omterm{def:b2:isotope-geochemistry:standards}{VSMOW}. وللكالسيت $^{18}\mathrm{O}/^{16}\mathrm{O} =
\num{2.0672e-3}\times(1 - 0.0020) = \num{2.0631e-3}$، ومنه $\delta^{18}\mathrm{O} =
(\num{2.0631e-3}/\num{2.0052e-3} - 1)\times1000 = +28.9$ بالجزء في الألف (\omterm{def:b2:isotope-geochemistry:standards}{VSMOW}).
\end{solution}

\begin{solution}{exo:b2:isotope-geochemistry:3}
عند \qty{1.0}{Ga}: $e^{\lambda t} - 1 = 0.006545$، و$0.512638 - 0.1966\times0.006545 =
0.511351$. وعند \qty{4.567}{Ga}: $e^{\lambda t}-1 = 0.030243$، و$0.512638 -
0.1966\times0.030243 = 0.506692$.
\end{solution}

\begin{solution}{exo:b2:isotope-geochemistry:4}
$\delta\mathrm{D} = 8\times(-18.0) + 10 = -134$ بالجزء في الألف. وللبحيرة، $d = -42 -
8\times(-4.0) = -10$ بالجزء في الألف، أي أقل من قيمة الخط 10: فالماء يقع تحت \omterm{prop:b2:isotope-geochemistry:mwl}{خط الماء
الجوي}، كما تقع المياه السطحية المتبخّرة (فقد غادرت النظائر الخفيفة أولًا).
\end{solution}

\begin{solution}{exo:b2:isotope-geochemistry:5}
$Q = 10^4\times\num{6.524e-12}\times0.1966/0.512638 = \num{2.502e-8}$ لكل سنة، أي
\qty{25.0}{\per\giga\year}. وللغرانيت $\mathrm{Sm}/\mathrm{Nd} = 6.0/30 = 0.200$،
مقابل $0.145/0.457 = 0.317$ لكوندريتات CI: $f = 0.200/0.317 - 1 = -0.37$. وعلى مدى
\qty{0.5}{Gyr}: $\Delta\varepsilon = 25.0\times(-0.37)\times0.5 = -4.6$ وحدة.
\end{solution}

\begin{solution}{exo:b2:isotope-geochemistry:6}
$\varepsilon \approx 3.5 + 25.0\times(-0.45)\times2.7 = 3.5 - 30.4 = -26.9$، مقابل
$-27.3$: فالتقريب أعلى مما ينبغي بمقدار 0.4 وحدة (أقل سلبيةً). ومعظم
الفرق، نحو 0.3 وحدة، يأتي من $e^{\lambda t} > 1$: فالأم، ومن ثمّ
العجز في النموّ نسبةً إلى CHUR، كانت أوفر في الماضي منها اليوم
(\cref{exo:b2:isotope-geochemistry:10})؛ والباقي يأتي من تقريب $f$ والقيمة الابتدائية.
\end{solution}

\begin{solution}{exo:b2:isotope-geochemistry:7}
مع كسر كتلي $x$ من الغرانيت، $\mathrm{Sr} = 320x + 130(1-x)$
و$R = (320x\times0.7200 + 130(1-x)\times0.7030)/\mathrm{Sr}$. ومن أجل $x = 0.05$: Sr
\qty{139.5}{ppm}، و$R = 0.70495$؛ و$x = 0.10$: 149.0 و0.70665؛ و$x = 0.20$: 168.0
و0.70948. ومن برهان قضية المزج، $R\,\mathrm{Sr}$ خطي في Sr:
$R = 0.731632 - 3.7221/\mathrm{Sr}$، أي خط مستقيم ميله $-3.722$ (ppm)
ومقطعه 0.7316 بدلالة $1/\mathrm{Sr}$. وللتحقّق عند $x = 0.05$: $0.731632 - 3.7221/139.5
= 0.70495$.
\end{solution}

\begin{solution}{exo:b2:isotope-geochemistry:8}
عند \qty{20}{\degreeCelsius}: $\delta^{18}\mathrm{O} = 1000/1.0098 - 1000 = -9.70$ بالجزء
في الألف، و$\delta\mathrm{D} = 1000/1.084 - 1000 = -77.5$ بالجزء في الألف، والميل $77.5/9.70 = 7.99$.
وعند \qty{0}{\degreeCelsius}: $-11.56$ و$-99.9$ بالجزء في الألف، والميل 8.64. وكلاهما قريب من
الميل 8 \omterm{prop:b2:isotope-geochemistry:mwl}{لخط الماء الجوي}: فالتجزّؤ التوازني للماء
بين السائل والبخار وحده يجعل نسبة إزاحتي الديوتيريوم والأكسجين 18
قريبة من 8.
\end{solution}

\begin{solution}{exo:b2:isotope-geochemistry:9}
تحوي القشرة $0.0051\times27/1.25 = 0.110$، أي 11 في المئة من النيوديميوم. وموازنة
الكتلة: $0.110\times(-15) + 0.890\,\varepsilon_m = 0$، ومنه $\varepsilon_m = 1.65/0.890 =
+1.9$. ومع خزّان مستنزَف يحوي حصّة $x$ عند $+8.5$ والباقي عند 0:
$0.110\times(-15) + 8.5x = 0$، $x = 1.65/8.5 = 0.19$: أي نحو خُمس
النيوديميوم.
\end{solution}

\begin{solution}{exo:b2:isotope-geochemistry:10}
معدل النموّ هو $\lambda p\,e^{\lambda t}$، مع $e^{\lambda t} \approx 1 + \lambda t$؛
وبأخذ متوسطه على الفترة من الانفصال $T$ إلى الحاضر يكون العامل
$1 + \lambda T/2$، فالحدود المهملة حجمها النسبي $\lambda T/2$. وعلى مدى
عمر الأرض، $\num{6.524e-12}\times\num{4.567e9}/2 = 0.015$، أي 1.5 في المئة. ومن أجل
قشرة \cref{exo:b2:isotope-geochemistry:6}، $\lambda T/2 = 0.0088$، والخطأ $0.0088\times30.4 = 0.27$
وحدة. ويضيف التغيّر البطيء لنسبة CHUR التي تُقسم عليها $\varepsilon$
حدًّا آخر أصغر، من رتبة $(p/R)\lambda T \times \varepsilon_0$، أي هنا 0.02 وحدة.
\end{solution}

\begin{solution}{exo:b2:isotope-geochemistry:11}
$K = 13.3/5.0 = 2.66$. ومع $s = 0.5$: $s' = 2.66\times0.5/(1 + 1.66\times0.5) =
1.33/1.83 = 0.73$. ويكون الفرق $s' - s$ أكبر ما يمكن حيث
$\mathrm{d}s'/\mathrm{d}s = K/(1+(K-1)s)^2 = 1$، أي عند $s = (\sqrt K - 1)/(K-1) =
0.631/1.66 = 0.38$، حيث $s' = 0.62$: فالفرق 0.24.
\end{solution}

\begin{solution}{exo:b2:isotope-geochemistry:12}
التوليد: من $2.3$ إلى $2.6\times\qty{7.2e9}{km^3}/\qty{4.5e9}{yr} = 3.7$ إلى
\qty{4.2}{km^3} في السنة. والهدم: من $1.3$ إلى $1.6\times7.2/4.5 = 2.1$ إلى
\qty{2.6}{km^3} في السنة. والمعدلان الحاليان مختلفان: فالإضافات (2.7) أدنى من متوسط
التوليد في النموذج، والإزالة (3.2) أعلى من متوسط الهدم فيه. ففي النموذج نمت القشرة
مبكرًا، حين كان التوليد أسرع، والحاضر زمن توازن أو فقد
طفيف، كما تقول موازنة مناطق الاندساس.
\end{solution}

\begin{solution}{pb:b2:isotope-geochemistry:1}
\textbf{1.} الوشاح: $(0.51307/0.512638 - 1)\times10^4 = +8.4$. الرسوب:
$(0.51141/0.512638 - 1)\times10^4 = -24.0$.

\textbf{2.} كورا $+2.8$، فيل $-4.3$، آشا $+4.7$.

\textbf{3.} مع $^{87}\mathrm{Sr}/^{86}\mathrm{Sr}$ أفقيًا و$\varepsilon_{\mathrm{Nd}}$ رأسيًا:
الوشاح (0.70282، $+8.4$)، وكورا (0.70410، $+2.8$)، وفيل (0.70618، $-4.3$)، والرسوب
(0.72000، $-24.0$)، وآشا (0.70290، $+4.7$). وفيل وحدها، بقيمة
$\varepsilon_{\mathrm{Nd}}$ السالبة، تقع في الجانب المُثرى من CHUR.

\textbf{4.} تحت المصفوفة: فعند نسبة السترونشيوم التي للوشاح المستنزَف تكون للمصفوفة
$\varepsilon_{\mathrm{Nd}}$ قرب $+8$؛ أمّا آشا فلها $+4.7$. وهي تشبه
مجموعة سانت هيلينا من نمط HIMU (مانغايا، توبواي، سانت هيلينا).

\textbf{5.} يغيّر الانصهار نسب الأم إلى الوليد لكنه لا يغيّر، على نحو قابل للقياس،
النسب النظائرية للسترونشيوم والنيوديميوم (القضية الخاصة بالانصهار)؛ والحمم
فتيّة، فلم تنمُ نسبها منذ الانبثاق (السؤال 6).

\textbf{6.} $\Delta R = 0.12\,(e^{\lambda\cdot10^6} - 1) = 0.12\times\num{6.52e-6} =
\num{7.8e-7}$، أي $\num{7.8e-7}/0.512638\times10^4 = 0.015$ وحدة إبسيلون: وهذا أدنى بكثير من
الفروق بين الجزر، وأدنى من دقة القياس.

\textbf{7.} $\mathrm{Nd} = 40f + 0.75(1-f)$،
و$R = \bigl(40f\times0.51141 + 0.75(1-f)\times0.51307\bigr)/\mathrm{Nd}$.

\textbf{8.} $f = 0.5$ في المئة: Nd \qty{0.9463}{ppm}، $\varepsilon = +1.6$؛ و1 في المئة:
1.1425، $-2.9$؛ و2 في المئة: 1.535، $-8.4$؛ و5 في المئة: 2.7125، $-15.4$.

\textbf{9.} Sr هو 10.95 و11.90 و13.80 و\qty{19.50}{ppm}؛ و$^{87}\mathrm{Sr}/^{86}\mathrm{Sr} =
0.70439$ و0.70571 و0.70780 و0.71163.

\textbf{10.} $R\,\mathrm{Sr} = 10\times0.70282 + f(200\times0.72000 - 10\times0.70282)$
و$\mathrm{Sr} = 10 + 190f$؛ وبحذف $f$، $R\,\mathrm{Sr} = 7.0282 +
(\mathrm{Sr}-10)\times136.9718/190$، ومنه $R = 0.720904 - 0.18084/\mathrm{Sr}$: أي خط
مستقيم ميله $-0.1808$ (ppm) بدلالة $1/\mathrm{Sr}$. وللتحقّق عند $f = 1$ في المئة:
$0.720904 - 0.18084/11.90 = 0.70571$.

\textbf{11.} $K = (10/0.75)/(200/40) = 13.33/5.0 = 2.67$. فالوشاح أغنى نسبيًا
بالسترونشيوم، والرسوب أغنى بالنيوديميوم، فتتحرّك نسبة Nd نحو
الرسوب أسرع: فقليل من الرسوب يزيح $\varepsilon_{\mathrm{Nd}}$ مسافة طويلة قبل أن
تتحرّك $^{87}\mathrm{Sr}/^{86}\mathrm{Sr}$ كثيرًا.

\textbf{12.} أغنى بالنيوديميوم 53 مرة ($40/0.75$)، وبالسترونشيوم 20 مرة. فعند $f = 1$ في المئة
يقدّم الرسوب $0.40/1.1425 = 35$ في المئة من النيوديميوم، وهذا يزيح
$\varepsilon_{\mathrm{Nd}}$ بمقدار $0.35\times(8.4 + 24.0) = 11.3$ وحدة، لكنه لا يقدّم إلا $2.0/11.9 =
17$ في المئة من السترونشيوم.

\textbf{13.} $0.01\times40/1.1425 = 0.350$: أي 35 في المئة.

\textbf{14.} من $R\,(40f + 0.75(1-f)) = 40fR_S + 0.75(1-f)R_D$:
$f = 0.75(R - R_D)/\bigl(40(R_S - R) + 0.75(R - R_D)\bigr)$، والشيء نفسه مع 200
و10 للسترونشيوم.

\textbf{15.} $f = 0.75\times(-0.00029)/\bigl(40\times(-0.00137) + 0.75\times(-0.00029)\bigr)
= \num{-2.175e-4}/\num{-5.502e-2} = 0.0040$: أي 0.40 في المئة.

\textbf{16.} $f = 10\times0.00128/\bigl(200\times0.01590 + 10\times0.00128\bigr) =
0.0128/3.1928 = 0.0040$: أي 0.40 في المئة؛ والتقديران متّفقان.

\textbf{17.} النيوديميوم: $0.75\times(-0.00065)/\bigl(40\times(-0.00101) +
0.75\times(-0.00065)\bigr) = \num{-4.875e-4}/\num{-4.089e-2} = 0.0119$؛ والسترونشيوم:
$10\times0.00336/\bigl(200\times0.01382 + 10\times0.00336\bigr) = 0.0336/2.7976 = 0.0120$.
وكلاهما يعطي 1.2 في المئة.

\textbf{18.} $f = 0.75\times(-0.00019)/\bigl(40\times(-0.00147) + 0.75\times(-0.00019)\bigr)
= 0.0024$، وهذا يتنبّأ بتركيز Sr قدره \qty{10.46}{ppm} وبنسبة $^{87}\mathrm{Sr}/^{86}\mathrm{Sr} =
0.70361$، مقابل 0.70290 مقيسة. فآشا ليست على منحنى المزج بين الوشاح والرسوب:
فمصدرها يحوي مكوّنًا آخر.

\textbf{19.} قشرة محيطية قديمة معاد تدويرها، من نمط HIMU: تحتاج نظائر Nd فيها إلى نسبة Sm/Nd
أدنى قليلًا من نسبة \omterm{def:b2:isotope-geochemistry:reservoirs}{الوشاح المستنزَف}، ونظائر Sr إلى نسبة Rb/Sr منخفضة،
ورصاصها (غير المقيس هنا) إلى نسبة U/Pb عالية، كل ذلك على مدى مليار إلى ملياري سنة.

\textbf{20.} $e^{\lambda_{147}\cdot2\times10^9} - 1 = 0.013133$. الرسوب: $0.51141 -
0.110\times0.013133 = 0.509965$؛ والخزّان CHUR: $0.512638 - 0.1966\times0.013133 = 0.510056$؛
$\varepsilon_{\mathrm{Nd}}(2\,\mathrm{Ga}) = (0.509965/0.510056 - 1)\times10^4 = -1.8$.

\textbf{21.} $f = 0.110/0.1966 - 1 = -0.440$؛ ومنذ الاندساس تغيّرت $\varepsilon$ بمقدار
$25.0\times(-0.440)\times2.0 = -22.0$، ومنه $\varepsilon(2\,\mathrm{Ga}) \approx -24.0 + 22.0
= -2.0$، في حدود 0.2 وحدة من القيمة الدقيقة.

\textbf{22.} $e^{\lambda_{87}\cdot2\times10^9} - 1 = 0.028338$؛
$0.72000 - 0.50\times0.028338 = 0.70583$.

\textbf{23.} كانت للرسوب عند الاندساس $\varepsilon_{\mathrm{Nd}} = -1.8$، قريبة من
قيمة $-4.7$ لقشرة لم يكن عمرها حينئذ إلا 0.7 مليار سنة، وأعلى بكثير من
القيم الشديدة السلبية للقشرة القديمة: فهي تلائم تعرية قشرة فتيّة، ربما
مختلطة ببعض المادة المشتقة من الوشاح. ونسبة $^{87}\mathrm{Sr}/^{86}\mathrm{Sr}$ فيه البالغة
0.7058 في ذلك الوقت هي أيضًا نسبة صخر قشري فتيّ.

\textbf{24.} يحوي مصدر فيل، الجزيرة الأكثر إثراءً،
\textbf{\qty[text-series-to-math=true]{1.2}{\percent}} من الرسوب المعاد تدويره.
\end{solution}
